El apartado ``Requisitos de las acreditaciones internacionales EUR-ACE y ABET'' de la Normativa de TFT de la ETSIT-UPM establece que ``\textit{La memoria del TFT del GITST, GIB y MUIT, y en general la de aquellas titulaciones que hayan obtenido o para las que se desee solicitar una acreditación internacional EUR-ACE o ABET, … debe incluir un presupuesto económico}''.

El presupuesto de este trabajo se recoge en la tabla~\ref{presupuesto-economico}. Los criterios
aplicados para elaborarlo son los siguientes:

\begin{itemize}
    \item \textbf{Mano de obra.} El trabajo se desarrolló a lo largo de veinte \textit{sprints} de
    dos semanas (de octubre a julio), diez meses en total. El coste imputado es el real: la beca de
    colaboración con la que se retribuyó al autor, de 500~\euro\ mensuales.

    \item \textbf{Recursos materiales amortizables.} Equipamiento e instrumentación del laboratorio
    y licencias de las herramientas de diseño, imputados en proporción a los meses de uso efectivo
    sobre su periodo de amortización. Las licencias de Vivado ML Enterprise y Altium Designer se
    amortizan a un año por tratarse de suscripciones anuales.

    \item \textbf{Material fungible.} Componentes electrónicos y fabricación de las placas de
    circuito impreso. Las cifras de componentes proceden directamente de las listas de materiales
    de los tres diseños, para una tirada de dos unidades de cada placa, tomando el precio del
    proveedor primario. Los costes de fabricación (placa de cuatro capas más \textit{stencil} en
    PCBWay~\cite{pcbway_caps}) son estimados.

    \item \textbf{Gastos generales y beneficio industrial.} Se aplican los porcentajes habituales
    del 15\,\% sobre el coste directo y del 6\,\% sobre la suma de costes directos e indirectos.
    El IVA aplicable es del 21\,\%.
\end{itemize}

\begin{table}[H]
    \centering
    \resizebox{\textwidth}{!}{
\begin{tabular}{lllllc}
\cline{4-6}

\multicolumn{3}{l|}{\textbf{COSTE DE MANO DE OBRA (coste directo)}} &
\multicolumn{1}{c|}{\textbf{Meses}} & \multicolumn{1}{c|}{\textbf{Importe mensual}} &
\multicolumn{1}{c|}{\textbf{Total}} \\ \cline{4-6}
\hline

\multicolumn{3}{|l|}{Becario de colaboración (desarrollo VHDL, firmware y PCB)} &
\multicolumn{1}{c|}{10} &
\multicolumn{1}{c|}{500,00 \euro} &
\multicolumn{1}{c|}{\textbf{5.000,00 \euro}} \\
\hline

\multicolumn{5}{|l|}{\textbf{COSTE TOTAL DE MANO DE OBRA}} &
\multicolumn{1}{c|}{\textbf{5.000,00 \euro}} \\
\hline

& & & & &
\multicolumn{1}{l}{} \\ \cline{3-6}

\multicolumn{2}{l|}{\textbf{COSTE DE RECURSOS MATERIALES (coste directo)}} &
\multicolumn{1}{c|}{\textbf{Precio de compra}} &
\multicolumn{1}{c|}{\textbf{Uso en meses}} &
\multicolumn{1}{c|}{\textbf{Amortización (en años)}} &
\multicolumn{1}{c|}{\textbf{Total}} \\
\hline

\multicolumn{2}{|l|}{Placa de desarrollo Xilinx ZCU102} &
\multicolumn{1}{c|}{3.500,00 \euro} &
\multicolumn{1}{c|}{10} &
\multicolumn{1}{c|}{5} &
\multicolumn{1}{c|}{583,33 \euro} \\
\hline

\multicolumn{2}{|l|}{Estación de trabajo (síntesis y diseño de PCB)} &
\multicolumn{1}{c|}{1.800,00 \euro} &
\multicolumn{1}{c|}{10} &
\multicolumn{1}{c|}{5} &
\multicolumn{1}{c|}{300,00 \euro} \\
\hline

\multicolumn{2}{|l|}{Osciloscopio digital de 4 canales} &
\multicolumn{1}{c|}{2.500,00 \euro} &
\multicolumn{1}{c|}{4} &
\multicolumn{1}{c|}{8} &
\multicolumn{1}{c|}{104,17 \euro} \\
\hline

\multicolumn{2}{|l|}{Analizador lógico USB} &
\multicolumn{1}{c|}{400,00 \euro} &
\multicolumn{1}{c|}{4} &
\multicolumn{1}{c|}{5} &
\multicolumn{1}{c|}{26,67 \euro} \\
\hline

\multicolumn{2}{|l|}{Fuente de alimentación de laboratorio} &
\multicolumn{1}{c|}{600,00 \euro} &
\multicolumn{1}{c|}{4} &
\multicolumn{1}{c|}{8} &
\multicolumn{1}{c|}{25,00 \euro} \\
\hline

\multicolumn{2}{|l|}{Estación de soldadura y horno de reflujo} &
\multicolumn{1}{c|}{1.200,00 \euro} &
\multicolumn{1}{c|}{3} &
\multicolumn{1}{c|}{8} &
\multicolumn{1}{c|}{37,50 \euro} \\
\hline

\multicolumn{2}{|l|}{Multímetro de mesa} &
\multicolumn{1}{c|}{300,00 \euro} &
\multicolumn{1}{c|}{3} &
\multicolumn{1}{c|}{8} &
\multicolumn{1}{c|}{9,38 \euro} \\
\hline

\multicolumn{2}{|l|}{Licencia Vivado ML Enterprise (anual)} &
\multicolumn{1}{c|}{2.000,00 \euro} &
\multicolumn{1}{c|}{10} &
\multicolumn{1}{c|}{1} &
\multicolumn{1}{c|}{1.666,67 \euro} \\
\hline

\multicolumn{2}{|l|}{Licencia Altium Designer (anual)} &
\multicolumn{1}{c|}{3.000,00 \euro} &
\multicolumn{1}{c|}{5} &
\multicolumn{1}{c|}{1} &
\multicolumn{1}{c|}{1.250,00 \euro} \\
\hline

\multicolumn{5}{|l|}{\textbf{COSTE TOTAL DE RECURSOS MATERIALES}} &
\multicolumn{1}{c|}{\textbf{4.002,72 \euro}} \\
\hline

& & & & &
\multicolumn{1}{l}{} \\
\hline

\multicolumn{5}{|l|}{\textbf{COSTE DIRECTO (CD) = mano de obra + recursos materiales}} &
\multicolumn{1}{c|}{\textbf{9.002,72 \euro}} \\
\hline

& & & & &
\multicolumn{1}{l}{} \\
\hline

\multicolumn{1}{|l|}{\textbf{GASTOS GENERALES (costes indirectos)}} &
\multicolumn{1}{c|}{15 \%} &
\multicolumn{3}{c|}{sobre CD} &
\multicolumn{1}{c|}{\textbf{1.350,41 \euro}} \\
\hline

\multicolumn{1}{|l|}{\textbf{BENEFICIO INDUSTRIAL}} &
\multicolumn{1}{c|}{6 \%} &
\multicolumn{3}{c|}{sobre CD + CI} &
\multicolumn{1}{c|}{\textbf{621,19 \euro}} \\
\hline

& & & & &
\multicolumn{1}{l}{} \\

\multicolumn{6}{l}{\textbf{MATERIAL FUNGIBLE}} \\
\hline

\multicolumn{5}{|l|}{Componentes de la placa de comunicación serie (2 unidades)} &
\multicolumn{1}{c|}{\textbf{424,13 \euro}} \\
\hline

\multicolumn{5}{|l|}{Componentes de la placa CDHS (2 unidades)} &
\multicolumn{1}{c|}{\textbf{116,77 \euro}} \\
\hline

\multicolumn{5}{|l|}{Componentes de la placa AOCS (2 unidades)} &
\multicolumn{1}{c|}{\textbf{302,02 \euro}} \\
\hline

\multicolumn{5}{|l|}{Fabricación de las tres placas con \textit{stencil} (2 unidades de cada una)} &
\multicolumn{1}{c|}{\textbf{420,00 \euro}} \\
\hline

\multicolumn{5}{|l|}{Cableado, conectores, tarjeta microSD y consumibles de soldadura} &
\multicolumn{1}{c|}{\textbf{180,00 \euro}} \\
\hline

\multicolumn{5}{|l|}{\textbf{COSTE TOTAL DE MATERIAL FUNGIBLE}} &
\multicolumn{1}{c|}{\textbf{1.442,92 \euro}} \\
\hline

& & & & &
\multicolumn{1}{l}{} \\
\hline

\multicolumn{5}{|l|}{\textbf{SUBTOTAL PRESUPUESTO}} &
\multicolumn{1}{c|}{\textbf{12.417,24 \euro}} \\
\hline

\multicolumn{4}{|l|}{\textbf{IVA APLICABLE}} &
\multicolumn{1}{c|}{21 \%} &
\multicolumn{1}{c|}{\textbf{2.607,62 \euro}} \\
\hline

& & & & &
\multicolumn{1}{l}{} \\
\hline

\multicolumn{5}{|l|}{\textbf{TOTAL PRESUPUESTO}} &
\multicolumn{1}{c|}{\textbf{15.024,86 \euro}} \\
\hline

\end{tabular}}
    \caption{Presupuesto económico.}
    \label{presupuesto-economico}

\end{table}

El coste total asciende a \textbf{quince mil veinticuatro euros con ochenta y seis céntimos}
(15.024,86~\euro), IVA incluido. La mano de obra es la mayor partida individual, con un 40\,\% del
subtotal, seguida de los recursos materiales amortizados con un 32\,\%, dentro de los cuales las
licencias de las herramientas de diseño pesan más que el propio equipamiento; el hardware
físicamente fabricado representa un 10\,\% del presupuesto.

El coste de mano de obra imputado es el realmente incurrido, correspondiente a una beca de
colaboración, y no el que tendría el mismo trabajo contratado a precio de mercado. Con
una tarifa de ingeniero junior de 30~\euro/hora sobre una dedicación equivalente de 800~horas, la
partida de mano de obra pasaría de 5.000~\euro\ a 24.000~\euro\ y el total del presupuesto se
situaría en torno a los 43.000~\euro. Esa segunda cifra es la referencia adecuada para estimar lo
que costaría reproducir este desarrollo en un entorno industrial.
